\documentclass[../../../../SoftwareRequirementsSpecification.tex]{subfiles}
\begin{document}
% Richiede le macro definite in src/macros/useCaseMacros.tex
% (incluse nel preambolo di SoftwareRequirementsSpecification.tex)

\ucstart
\begin{xltabular}{\textwidth}{|l|X|}
  \hline
  \ucid{PT-10} \\ \hline
  \ucname{Modifica Scheda} \\ \hline
  \ucactor{PT} \\ \hline
  \ucdescription{Il PT modifica il nome e la descrizione di una
    scheda che ha creato.} \\ \hline
  \ucprecond{Il PT deve essere autenticato. Deve esistere almeno una
    scheda da modificare.} \\ \hline
  \ucpostcond{Il sistema aggiorna il nome e/o la descrizione della
    scheda.} \\ \hline
  \ucmainflow{
    \item Il PT apre la pagina delle proprie schede.
    \item Il sistema mostra l'elenco delle schede create dal PT.
    \item Il PT preme il pulsante "Modifica" su una scheda
      dell'elenco.
    \item Il sistema mostra il form di modifica precompilato con il
      nome e la descrizione correnti della scheda.
    \item Il PT modifica il nome e/o la descrizione.
    \item Il PT preme il pulsante "Salva Modifiche".
    \item Il sistema verifica che, se il nome è cambiato, non sia
      già usato da un'altra scheda del PT.
    \item Il sistema salva le modifiche.
    \item Il sistema mostra un messaggio di conferma.
  } \\ \hline
  \ucaltflow{
    \textbf{7a. Nome già presente} \newline
    - Il sistema mostra un messaggio di errore ("Esiste già una
    scheda con questo nome"). \newline
    - Il flusso riprende dal passo 4.
  } \\ \hline
  \ucnote{
    Questo caso d'uso modifica solo il nome e la descrizione della
    scheda. Le sessioni di allenamento si modificano con il caso
    d'uso PT-04, indipendentemente dal fatto che la scheda abbia
    assegnazioni attive o meno: il nome e la descrizione non hanno
    effetto su ciò che gli atleti stanno seguendo, quindi non c'è
    bisogno di restrizioni come quelle di PT-04.
  } \\ \hline
\end{xltabular}
\end{document}
